\subsection{Learned Representations and Model Explainability}

To assess what the CRPS ensemble U-Net learns from the 19 ASCAT predictors and
stochastic latent input, Fig.~\ref{fig:explainability} combines bottleneck
representation analysis, predictor-group permutation with uncertainty,
spatial attribution, and ensemble coherence.
Bottleneck features (fixed latent draw) reduced by PCA separate high
exceedance-probability pixels from background and organize ambiguity count
$N_{\mathrm{amb}}$ along distinct directions of the embedding
(Fig.~\ref{fig:explainability}a,b).
Along-track permutation of predictor groups shows that
$N_{\mathrm{amb}}$ and ambiguity winds/SOS contribute most to the exceedance
Brier score at $\tau^*=1.2\times10^{-4}$, exceeding smooth wind context, with
uncertainties small relative to group differences
(Fig.~\ref{fig:explainability}c); corresponding afCRPS permutation numbers are
retained in the Supporting Information.
Integrated gradients on a representative orbit attribute high $P(S>\tau^*)$ to
coherent gust-front-scale structure rather than isolated retrieval artifacts
(Fig.~\ref{fig:explainability}d).
Member-difference correlation lengths of ensemble deviations concentrate near
5--6~pixels, at the latent-grid spacing rather than at long front-scale modes
(Fig.~\ref{fig:explainability}e).
Additional PCA colorings, the full IG case set, and latent $z$-traversal
examples are provided in the Supporting Information
(Figs.~\ref{fig:si-explain-pca}--\ref{fig:si-explain-ig-share}).
